% Section 5, moved as is from old Sections 4.6-4.7 (loop0009 item 01). Sources:
% paper2/solver_fix.md (cost, re-certification), paper2/certificates.md (proof
% object), paper2/dataset.md (the dataset). Tables from tables/ (make_tables.py).
\section{Computational results and a dataset}
\label{sec:results}

\subsection{The revised algorithm in practice}
\label{sec:practice}

The revised search changes two tests. The definite move fires on the first
playable $q$ that passes the published test and then the matching of
Theorem~\ref{thm:matching}; a $q$ that fails the matching is passed over. The
better move cites $q$ for $r$ only if the same matching succeeds at
$\cl(S\cup\{r\})$. The matching joins at most $\close-1$ customers to $\open$
stacks and stops at $\open-1$ edges. Since 1 October 2026 these are the default
in both of our solvers, the MOSP customer search and the graph pathwidth solver,
in C and in Python; the published rules remain behind a flag.

\begin{table}[t]
\centering\small
\caption{The cost of the repair: nodes under the published and the repaired
rules, in paired runs of the same decision in one worker. Totals over the pairs
that finished under both settings; a censored pair says nothing about cost.
MOSP rows refute the optimum minus one; the graph row is whole descents at
120\,s each.}
\label{tab:cost}
\resizebox{\linewidth}{!}{% Generated by python -m paper2.latex.make_tables. Source: paper2/data/solver_fix_cost_{mosp40,cs,cs125,pw}.csv.
% Do not edit by hand.
\begin{tabular}{@{}llrrrrr@{}}
\toprule
Where & Vertices & Pairs finished & Nodes, published & Nodes, repaired & Ratio & Worst pair \\
\midrule
MOSP corpus, \texttt{default} & 9--40 & 6,135 of 6,135 & 228,147 & 228,154 & 1.00003 & 1.009 \\
MOSP corpus, \texttt{csearch} & 9--40 & 6,135 of 6,135 & 202,962 & 202,971 & 1.00004 & 1.040 \\
Chu \& Stuckey classes & 50--100 & 118 of 125 & 1,358,096,671 & 1,362,205,210 & 1.00303 & 1.013 \\
Chu \& Stuckey classes, cap $2\times10^8$ nodes & 125 & 11 of 23 & 90,969,508 & 90,961,419 & 0.99991 & 1.000 \\
Graph pathwidth, whole descents & 4--2{,}916 & 731 of 880 & 2,977,849,807 & 3,038,670,717 & 1.020 & -- \\
\bottomrule
\end{tabular}
}
\end{table}

Table~\ref{tab:cost} gives the cost in nodes. The repair changes tree sizes by
at most a few percent and never changed an answer. Per node it costs about 2\% in
the MOSP C engine, measured on six hard instances at 75--125 customers. The
published test passes and the matching then fails on 0.07--0.27\% of
definite-move candidates, a rate that rises with size.

\paragraph{Re-certification.} Of the 6,374 certified optima in our MOSP corpus,
6,259 do not depend on the customer search: they rest on an exact subset-lattice
computation, a checked DRAT refutation, an earlier SAT search, or a lower bound
equal to the value. The other 115 rested on the customer search under the
published rules. The repaired solver has re-refuted the value minus one on 112
of them, and no value changed. Table~\ref{tab:split} gives the seven hardest,
125-customer instances at densities 2 and 4, which needed a split of the root
over many cores; three are still running at the time of writing, and keep their
values as verified upper bounds until they finish. The split is proved
sound~\leanref{split}.

\begin{table}[t]
\centering\small
\caption{Re-certification of the hardest corpus values under the repaired rules,
by a split of the root into independent tasks. \emph{$k$}: the value minus one,
which must be refuted. For an instance in flight, the nodes and core-hours are
those of the finished tasks so far.}
\label{tab:split}
% Generated by python -m paper2.latex.make_tables. Source: paper2/data/solver_fix_split_{tasks,results}.csv, read 2026-10-03 09:01.
% Do not edit by hand.
\begin{tabular}{@{}lrrlrrr@{}}
\toprule
Instance & Value & $k$ & State & Tasks refuted & Nodes & Core-hours \\
\midrule
\texttt{Random-125-125-4-1\_0} & 57 & 56 & refuted & 669 & $4.97\times10^{10}$ & 18.1 \\
\texttt{Random-125-125-2-4\_0} & 24 & 23 & refuted & 1,686 & $4.98\times10^{11}$ & 154.1 \\
\texttt{Random-125-125-4-5\_0} & 46 & 45 & refuted & 616 & $7.96\times10^{10}$ & 26.1 \\
\texttt{Random-125-125-4-2\_0} & 57 & 56 & refuted & 842 & $7.64\times10^{10}$ & 30.2 \\
\texttt{Random-125-125-2-1\_0} & 24 & 23 & in flight & 197 & $8.24\times10^{10}$ & 51.0 \\
\texttt{Random-125-125-2-5\_0} & 20 & 19 & in flight & 168 & $6.97\times10^{10}$ & 45.6 \\
\texttt{Random-125-125-4-4\_0} & 51 & 50 & in flight & 745 & $1.28\times10^{11}$ & 67.1 \\
\bottomrule
\end{tabular}
% read 2026-10-03 09:01; last task stamp 2026-10-03T09:00:59

\end{table}

\paragraph{Certificates.} A refutation by the search is a claim about a tree of
$10^{11}$ nodes at 125 customers, and nothing in Theorem~\ref{thm:main} lets a
reader check that the code ran the search it models. So the search emits a
certificate: the tree, every pruning step with its rule and witness, and for
each definite and better step the matching of Theorem~\ref{thm:matching}. An
independent checker, which shares no code with the search, recomputes every
premise from the matrix, verifies each matching (distinct customers, distinct
stacks, each stack newly opened by its customer, at least $\open-1$ edges),
requires every chain of covers to end at a child or an old move with no cycle,
and replays the tree. Each accepted link is a Lean theorem; that an acyclic
covering by such links is sound is a short paper argument, not yet formal.
Table~\ref{tab:certs} gives the results on the corpus at 9--75 customers.
Every refutation emitted within 60\,s verifies, 6,275 of 6,286 under the
\texttt{csearch} configuration, and none is rejected; the 35.5\,MB of
certificates check in 147\,s. On the graph of Counterexample~\ref{cex:definite}
rooted at $\{2\}$, the published rules emit an exhaustive tree that claims no
solution; a checker of the published premises accepts it, and ours rejects it,
even with a largest matching attached (one edge where two are needed).

\begin{table}[t]
\centering\small
\caption{Certificates for the repaired rules on the certified corpus at 9--75
customers, refuting the optimum minus one, 60\,s emission budget.
\emph{Unknown}: emission did not finish. Below 41 customers the
\texttt{default} rows match \texttt{csearch} and are omitted.}
\label{tab:certs}
\resizebox{\linewidth}{!}{% Generated by python -m paper2.latex.make_tables. Source: paper2/data/certificates/repaired.csv.
% Do not edit by hand.
\begin{tabular}{@{}llrrrrrrrrr@{}}
\toprule
Customers & Config. & Inst. & Verified & Rejected & Unknown & Definite & Better & Matching edges & gz bytes (median) & Check ms (median) \\
\midrule
9--10 & \texttt{csearch} & 1,614 & 1,614 & 0 & 0 & 307 & 92 & 32 & 332 & 0.04 \\
11--20 & \texttt{csearch} & 2,508 & 2,508 & 0 & 0 & 2,657 & 1,171 & 658 & 356 & 0.08 \\
21--30 & \texttt{csearch} & 1,816 & 1,816 & 0 & 0 & 14,796 & 6,148 & 5,675 & 414 & 0.24 \\
31--40 & \texttt{csearch} & 197 & 197 & 0 & 0 & 46,337 & 18,176 & 17,150 & 650 & 1.27 \\
41--50 & \texttt{csearch} & 120 & 120 & 0 & 0 & 530,424 & 120,302 & 152,597 & 12,045 & 33.24 \\
51--60 & \texttt{csearch} & 1 & 1 & 0 & 0 & 5,155 & 2,709 & 1,885 & 126,508 & 278.86 \\
61--75 & \texttt{csearch} & 30 & 19 & 0 & 11 & 103,095 & 31,404 & 37,261 & 26,516 & 143.70 \\
\midrule
41--50 & \texttt{default} & 120 & 120 & 0 & 0 & 748,952 & 0 & 191,780 & 12,665 & 37.01 \\
51--60 & \texttt{default} & 1 & 1 & 0 & 0 & 24,685 & 0 & 8,797 & 661,951 & 1198.43 \\
61--75 & \texttt{default} & 30 & 20 & 0 & 10 & 312,897 & 0 & 110,118 & 34,162 & 187.03 \\
\bottomrule
\end{tabular}
% total csearch: 6275 of 6286 verified, 0 rejected, 35.47 MB gz, 147.4 s, 3353245 nodes
% total default: 6276 of 6286 verified, 0 rejected, 54.77 MB gz, 242.3 s, 5361282 nodes
}
\end{table}

\subsection{A dataset with certified pathwidth}
\label{sec:dataset}

Section~\ref{sec:equiv} makes one certified pathwidth an exact answer to every
problem in Table~1's exact core: pathwidth, vertex separation and the vertex
separator game take the width $w$; MOSP, gate matrix layout with multiple
folding, one-dimensional logic, interval thickness, narrowness and node search
take $w+1$. We therefore built one dataset from every benchmark collection of
those problems that we could obtain, and the MOSP corpus. The unit is the
isomorphism class of the graph, the MOSP graph for a matrix instance: two
instances with isomorphic graphs have the same value under every problem, so the
dataset holds one record per class, with its members and their isomorphism maps.
Each record carries the graph, the width, its provenance (\emph{refutation},
\emph{bound} or \emph{upper bound}), a witness layout of exactly that width and,
for a bound, a minor certificate that replays.

\begin{table}[t]
\centering\small
\caption{The dataset by collection. \emph{Classes}: isomorphism classes within
the collection (nauty certificates; 78 graphs above 5,000 vertices by file
content only). \emph{Owned}: classes whose first member is in this collection,
so the column sums to the dataset's class count. The provenance columns count
owned classes with a witness layout.}
\label{tab:dataset}
\resizebox{\linewidth}{!}{% Generated by python -m paper2.latex.make_tables. Source: paper2/data/dataset/{index,classes}.csv.gz.
% Do not edit by hand.
\begin{tabular}{@{}lrrrrrrr@{}}
\toprule
Collection & Files & Classes & Owned & Refutation & Bound & Upper bound & No witnessed value \\
\midrule
MOSP: Challenge 2005 & 5,806 & 3,169 & 3,169 & 3,168 & 1 & 0 & 0 \\
MOSP: Challenge 2005 (second copy) & 46 & 46 & 0 & 0 & 0 & 0 & 0 \\
MOSP: Chu \& Stuckey & 200 & 200 & 200 & 198 & 0 & 2 & 0 \\
MOSP: Faggioli \& Bentivoglio & 300 & 287 & 274 & 274 & 0 & 0 & 0 \\
MOSP: SCOOP & 24 & 24 & 24 & 24 & 0 & 0 & 0 \\
VLSI gate matrix circuits & 11 & 11 & 11 & 5 & 5 & 1 & 0 \\
VSPLIB trees & 50 & 28 & 28 & 28 & 0 & 0 & 0 \\
VSPLIB grids & 50 & 50 & 50 & 9 & 0 & 19 & 22 \\
VSPLIB Harwell-Boeing & 73 & 72 & 72 & 32 & 7 & 33 & 0 \\
Small (Marti et al. 2008) & 84 & 84 & 84 & 32 & 52 & 0 & 0 \\
Rome graphs & 11,534 & 11,199 & 11,199 & 8,126 & 2,746 & 326 & 1 \\
TreewidthLIB colouring & 58 & 58 & 58 & 18 & 13 & 27 & 0 \\
freetdi named graphs & 150 & 146 & 143 & 97 & 22 & 23 & 1 \\
freetdi control-flow graphs & 1,817 & 1,070 & 1,064 & 398 & 661 & 4 & 1 \\
MOSP: Frinhani et al. 2018, large & 610 & 610 & 610 & 0 & 0 & 0 & 610 \\
PACE 2016 & 291 & 288 & 112 & 37 & 18 & 51 & 6 \\
PACE 2017 exact & 200 & 199 & 193 & 61 & 7 & 115 & 10 \\
PACE 2017 heuristic & 200 & 197 & 173 & 11 & 3 & 45 & 114 \\
PACE 2017 bonus & 100 & 100 & 100 & 23 & 0 & 77 & 0 \\
MOSP: Carvalho \& Soma 2015 & 150 & 150 & 150 & 11 & 0 & 139 & 0 \\
\midrule
All & 21,754 &  & 17,714 & 12,552 & 3,535 & 862 & 765 \\
\bottomrule
\end{tabular}
}
\end{table}

Table~\ref{tab:dataset} gives the counts. Twenty collections, 21,754 instance
files, reduce to 17,714 classes: the MOSP collections
\citep{SmithGent2005,FaggioliBentivoglio1998,ChuStuckey2009,CarvalhoSoma2015,Frinhani2018},
eleven VLSI gate matrix circuits \citep{OliveiraLorena2002}, the VSPLIB sets
\citep{Duarte2012} and the Small graphs \citep{MartiCamposPinana2008}, the Rome
graphs \citep{DiBattista1997}, and treewidth collections used as pathwidth
benchmarks \citep{TreewidthLIB,freetdi,PACE2016,PACE2017}. Of the classes, 16,087
carry a certified pathwidth with a witness layout and 862 a witnessed upper
bound. An independent checker, sharing no code with the solvers, re-reads all
20,987 member files with its own readers, maps each onto its record by
isomorphism, re-measures every witness and replays every bound certificate, and
reports no failure on the 16,949 records.

Against published values: all 84 Small graphs agree with the pathwidths of
\citet[Tables~4--5]{Mallach2018}, and ten of the eleven VLSI circuits are
certified at their published best-known track counts \citep[Table~I]{OliveiraLorena2002};
the eleventh, W4 with 202 nets, is open at 28 against 27. One caution about MOSP
benchmark files emerged: the same public repository ships the Carvalho and Soma
instances as customers by patterns and the Chu and Stuckey instances beside them
as patterns by customers, and only published densities and values tell them
apart. Read the wrong way round, an instance certifies at a different value.

What is left is priced. One more budget step for every open class that fits our
engine (1,024 vertices per component) costs at most 500 core-hours, and on the
Rome graphs each such step has closed about 40\% of what it attempted. The
largest MOSP instances, 400 to 1,000 customers, and the 154 classes above 1,024
vertices are beyond an exact search. Refutation-certified records rest on the
repaired search; for those within reach of the certificate emitter, the proof
object above is the route by which a third party can check them without our code.
